An oscillator turns DC into a periodic RF signal. Design means guaranteeing
start-up, setting the frequency, and minimizing phase noise.

\section{Two Views of the Same Thing}
The \emph{feedback} view (Barkhausen): loop gain $\ge 1$ with $0^\circ$ (mod
$360^\circ$) phase. The equivalent \emph{negative-resistance} view: the active
device presents a negative resistance $-R_d$ that cancels the resonator loss
$R_{res}$. Oscillation starts when $|R_d| > R_{res}$ and settles when, through
nonlinear compression, $|R_d| = R_{res}$ at the resonant frequency.

\section{Topologies}
The Colpitts, Hartley, and Clapp oscillators differ only in how the tank taps
feedback (a capacitive divider, an inductive divider, or a series-tuned
variant). For the best stability, a \emph{crystal} replaces the LC tank: its
enormous $Q$ ($10^4$--$10^6$) pins the frequency and slashes phase noise.
Dielectric-resonator (DRO) and YIG oscillators serve the microwave range.

\section{Design Flow}
Choose a resonator with the highest practical loaded $Q$; bias the device for
sufficient small-signal negative resistance (start-up margin $\sim$3$\times$);
set the tank for the target frequency; and buffer the output so the load does not
pull the frequency. Leeson's model (Chapter on Phase Noise) ties the achievable
phase noise to $Q$ and the device's flicker noise --- high $Q$ and a low-$1/f$
device win.
